\subsection{Choice and Zorn's Lemma}

\begin{definition} \label{def:choice-function} \hlt{Choice Function}\\
Let $\mathcal{F}$ be a family of nonempty sets, let $I$ be a set, and let $(X_i)_{i\in I}$ be an indexed family of nonempty sets.
\begin{enumerate}[label=(\roman*)]
\setlength{\itemsep}{0pt}
\item A \hlt{choice function} on $\mathcal{F}$ is a map $c:\mathcal{F}\to\bigcup\limits_{F\in\mathcal{F}}F$ satisfying $(\forall F\in\mathcal{F})\,c(F)\in F$.
\item An \hlt{indexed choice function} for $(X_i)_{i\in I}$ is a map $c:I\to\bigcup\limits_{i\in I}X_i$ satisfying $(\forall i\in I)\,c(i)\in X_i$.
\end{enumerate}
\end{definition}

\begin{principle} \label{princ:axiom-of-choice-analysis} \hlt{Choice Principle}\\
Let $I$ be a set and let $(X_i)_{i\in I}$ be an indexed family of nonempty sets. The axiom of choice asserts $(\exists c:I\to\bigcup\limits_{i\in I}X_i)(\forall i\in I)\,c(i)\in X_i$.\\
Equivalently, $\prod\limits_{i\in I}X_i\neq\varnothing$. An explicit $c$ proves existence directly; the principle supplies simultaneous choices without an explicit rule.
\end{principle}

\begin{principle} \label{princ:zorns-lemma} \hlt{Zorn's Lemma}\\
Let $(P,\leq)$ be a poset with $P\neq\varnothing$. Assume the following chain-upper-bound condition.\\
$(\forall C\subseteq P)\,[((\forall x,y\in C)\,(x\leq y\lor y\leq x))\Rightarrow(\exists u\in P)(\forall c\in C)\,c\leq u]$. Then $(\exists m\in P)(\forall x\in P)\,(m\leq x\Rightarrow x=m)$.
\end{principle}

\begin{definition} \label{def:inclusion-and-containment-orders-for-zorn} \hlt{Inclusion and Containment Orders}\\
Let $\mathcal{F}$ be a family of sets and let $\mathcal{A}\subseteq\mathcal{F}$.
\begin{enumerate}[label=(\roman*)]
\setlength{\itemsep}{0pt}
\item \hlt{Inclusion}: for $A,B\in\mathcal{F}$, define $A\leq B$ iff $A\subseteq B$. An upper bound of $\mathcal{A}$ is $U\in\mathcal{F}$ satisfying $(\forall A\in\mathcal{A})\,A\subseteq U$.\\
An element $M\in\mathcal{F}$ is maximal iff $(\forall B\in\mathcal{F})\,(M\subseteq B\Rightarrow B=M)$.
\item \hlt{Reverse inclusion}: for $A,B\in\mathcal{F}$, define $A\leq B$ iff $B\subseteq A$. An upper bound of $\mathcal{A}$ is $U\in\mathcal{F}$ satisfying $(\forall A\in\mathcal{A})\,U\subseteq A$.\\
An element $M\in\mathcal{F}$ is maximal iff $(\forall B\in\mathcal{F})\,(B\subseteq M\Rightarrow B=M)$.
\end{enumerate}
\end{definition}

\begin{proposition} \label{prop:inclusion-and-reverse-inclusion-partial-orders} \hlt{Inclusion Orders Are Partial Orders}\\
Let $\mathcal{F}$ be a family of sets. Inclusion and reverse inclusion are partial orders on $\mathcal{F}$.
\begin{proof}
Fix $A,B,C\in\mathcal{F}$. Since $A\subseteq A$, inclusion is reflexive. If $A\subseteq B$ and $B\subseteq C$, then $A\subseteq C$, so inclusion is transitive.\\
If $A\subseteq B$ and $B\subseteq A$, then $A=B$ by Assumption \ref{assump:set-extensionality}, so inclusion is antisymmetric. Hence inclusion is a partial order.\\
For reverse inclusion, $A\supseteq A$ gives reflexivity, $A\supseteq B\supseteq C$ gives $A\supseteq C$, and $A\supseteq B\supseteq A$ gives $A=B$ by Assumption \ref{assump:set-extensionality}.\\
Thus reverse inclusion is also a partial order.\\
\hlt{Structural Commentary}: Each order law follows from the corresponding containment law and set extensionality.\\
\hlt{Key Insight}: Reversing containment preserves reflexivity, transitivity, and antisymmetry.
\end{proof}
\end{proposition}

\begin{proofmethod} \label{meth:zorn-maximality-template} \hlt{Zorn Maximality Template}\\
Use this method when the target follows from a candidate admitting no proper admissible extension.\\
Let $P$ be the set of candidates, and let $\leq$ record admissible extension.
\begin{enumerate}[label=(\roman*)]
\setlength{\itemsep}{0pt}
\item Verify that $P\neq\varnothing$ and that $\leq$ is a partial order on $P$.
\item Fix a chain $C\subseteq P$, construct $u\in P$, and prove $(\forall p\in C)\,p\leq u$.
\item Obtain a maximal $m\in P$ by Principle \ref{princ:zorns-lemma}, and prove that maximality forces the target property.
\end{enumerate}
For $C\neq\varnothing$, common candidates are $\bigcup\limits_{A\in C}A$ under inclusion and $\bigcap\limits_{A\in C}A$ under reverse inclusion.\\
Under reverse inclusion, another object contained in every member of $C$ may be necessary.\\
If $C=\varnothing$, any element of $P$ is an upper bound. In every case, prove $u\in P$; this is the indispensable admissibility check.\\
\hlt{Key Idea}: A chain bound must preserve every admissibility condition defining the candidate set.
\end{proofmethod}
